\subsection{The \bsq{module} Comment}\label{sec:ModuleComment}
The documentation comment for a module is part of the \texttt{module-info.java} file for that module (refer to chapter \tqfullref{sec:ModuleDefinition}).

The \texttt{<module description>} is regular documentation comment, followed by the mdoule definition. It describes the module, its dependencies and what it provides. See the Javadoc\index{javadoc} tags \bdq{\nameref{sec:TagProvides}}\index{provides@"@provides} and \bdq{\nameref{sec:TagUses}}\index{uses@"@uses} for details.

The module documentation comment should not list the packages; this will be done automatically through the Javadoc\index{javadoc} tool during the generation process. Is should also not describe the packages; this is part of the package documentation.
 
If the project has just one module, the module comment can replace the overview comment (see \hyperref[sec:OverviewDoc]{above}).

%==============================================================================
% End of file!!
%==============================================================================

